\documentclass[11pt,a4paper]{article}

\def\courseTitle{Industrial Security (Bachelor)}
\def\labNumber{08}
\def\labTitle{Passive OT IDS with Zeek -- Tuning Rules Against Modbus Attacks}
\def\labSection{08}
\def\labTime{105--135 min}

% =====================================================================
% Shared preamble for lab sheets.
%
% Caller must \documentclass{article} first, then define:
%   \def\courseTitle{...}
%   \def\labNumber{NN}
%   \def\labTitle{...}
%   \def\labTime{e.g. "30--45 min"}
%   \def\labSection{e.g. "02"}
% Then % =====================================================================
% Shared preamble for lab sheets.
%
% Caller must \documentclass{article} first, then define:
%   \def\courseTitle{...}
%   \def\labNumber{NN}
%   \def\labTitle{...}
%   \def\labTime{e.g. "30--45 min"}
%   \def\labSection{e.g. "02"}
% Then % =====================================================================
% Shared preamble for lab sheets.
%
% Caller must \documentclass{article} first, then define:
%   \def\courseTitle{...}
%   \def\labNumber{NN}
%   \def\labTitle{...}
%   \def\labTime{e.g. "30--45 min"}
%   \def\labSection{e.g. "02"}
% Then \input{../../template/lab-preamble.tex}
% =====================================================================

\usepackage[utf8]{inputenc}
\usepackage[T1]{fontenc}
\usepackage[english]{babel}
\usepackage[a4paper, margin=2cm]{geometry}
\usepackage{graphicx}
\usepackage{listings}
\usepackage{xcolor}
\usepackage{colortbl}
\usepackage{tikz}
\usepackage{amsmath}
\usepackage{amssymb}
\usepackage{booktabs}
\usepackage{enumitem}
\usepackage{titlesec}
\usepackage{fancyhdr}
\usepackage{hyperref}
\usepackage{tcolorbox}
\tcbuselibrary{skins, breakable}
\usepackage{fontawesome5}

\input{../../../template/tikz-styles.tex}

\providecommand{\courseAuthor}{Industrial Security Lecture Series}
\providecommand{\courseInstitute}{Department of Computer Science}
\providecommand{\companyLogo}{logo}

\definecolor{slideAccent}{HTML}{00205B}
\definecolor{slideSecondary}{HTML}{00A0DC}

% --- Corporate branding overrides ------------------------------------
\IfFileExists{../../../branding.tex}{\input{../../../branding.tex}}{}

\colorlet{labAccent}{slideAccent}
\colorlet{labSec}{slideSecondary}
\definecolor{labCheck}{HTML}{2E7D32}
\definecolor{labWarn}{HTML}{B23A48}

\lstset{
  backgroundcolor=\color{black!4},
  basicstyle=\ttfamily\footnotesize,
  breaklines=true,
  frame=leftline,
  framerule=2pt,
  rulecolor=\color{labAccent},
  numbers=left,
  numberstyle=\tiny\color{black!50},
  showstringspaces=false,
  tabsize=2
}

\setlength{\tabcolsep}{4pt}

% Let TeX add a little extra inter-word stretch as a last resort before
% it reports an Overfull \hbox. Only kicks in on lines that would
% otherwise overflow (long \texttt paths, URLs, CIDRs); never loosens a
% line that already fits. Keeps prose/itemize within the margin.
\setlength{\emergencystretch}{3em}

\pagestyle{fancy}
\fancyhf{}
\renewcommand{\headrulewidth}{0.4pt}
\lhead{\small \courseTitle{} -- Lab \labNumber}
\rhead{\small Maps to Section \labSection}
\cfoot{\thepage}

\newtcolorbox{stepbox}[1]{
  enhanced, breakable, arc=2pt, boxrule=0pt,
  colback=labAccent!7, colframe=labAccent, leftrule=3pt,
  fonttitle=\bfseries\small, coltitle=white,
  colbacktitle=labAccent,
  title={\faTools\ Step #1}
}

\newtcolorbox{warnbox}{
  enhanced, breakable, arc=2pt, boxrule=0pt,
  colback=labWarn!7, colframe=labWarn, leftrule=4pt,
  fonttitle=\bfseries\small, coltitle=white, colbacktitle=labWarn,
  title={\faExclamationTriangle\ Caution}
}

\newtcolorbox{notebox}{
  enhanced, breakable, arc=2pt, boxrule=0pt,
  colback=labAccent!4, colframe=labAccent, leftrule=2pt,
  fonttitle=\bfseries\small, coltitle=white, colbacktitle=labAccent,
  title={\faInfoCircle\ Note}
}

\newtcolorbox{goalbox}{
  enhanced, breakable, arc=2pt, boxrule=0pt,
  colback=labCheck!7, colframe=labCheck, leftrule=3pt,
  fonttitle=\bfseries\small, coltitle=white, colbacktitle=labCheck,
  title={\faBullseye\ Goal}
}

\newtcolorbox{deliverbox}{
  enhanced, breakable, arc=2pt, boxrule=0pt,
  colback=labSec!10, colframe=labSec, leftrule=3pt,
  fonttitle=\bfseries\small, coltitle=white, colbacktitle=labSec,
  title={\faClipboardList\ Deliverable}
}

% --- Solution-sheet boxes (used only by lab-solutions.tex) ----------
\newtcolorbox{solutionbox}[1]{
  enhanced, breakable, arc=2pt, boxrule=0pt,
  colback=labCheck!7, colframe=labCheck, leftrule=3pt,
  fonttitle=\bfseries\small, coltitle=white, colbacktitle=labCheck,
  title={\faCheckCircle\ #1}
}

\newtcolorbox{rubricbox}{
  enhanced, breakable, arc=2pt, boxrule=0pt,
  colback=labSec!7, colframe=labSec, leftrule=3pt,
  fonttitle=\bfseries\small, coltitle=white, colbacktitle=labSec,
  title={\faBalanceScale\ Grading rubric}
}

\titleformat{\section}{\large\bfseries\color{labAccent}}{\thesection}{0.5em}{}

\newcommand{\labheader}{%
  \noindent
  \begin{tikzpicture}
    \fill[labAccent] (0,0) rectangle (\textwidth, -1.6cm);
    \node[anchor=north west, text=white, font=\bfseries\large] at (0.6cm, -0.4cm) {Lab \labNumber{} -- \labTitle};
    \node[anchor=south west, text=white, font=\small] at (0.6cm, -1.3cm) {\courseTitle{} \quad $\bullet$ \quad Maps to Section \labSection{} \quad $\bullet$ \quad \labTime};
    \node[anchor=north east, fill=labSec, text=white, font=\bfseries, inner sep=4pt, rounded corners=2pt]
      at ([xshift=-0.4cm, yshift=-0.4cm]\textwidth,0) {LAB \labNumber};
  \end{tikzpicture}
  \vspace{4mm}
}

% =====================================================================

\usepackage[utf8]{inputenc}
\usepackage[T1]{fontenc}
\usepackage[english]{babel}
\usepackage[a4paper, margin=2cm]{geometry}
\usepackage{graphicx}
\usepackage{listings}
\usepackage{xcolor}
\usepackage{colortbl}
\usepackage{tikz}
\usepackage{amsmath}
\usepackage{amssymb}
\usepackage{booktabs}
\usepackage{enumitem}
\usepackage{titlesec}
\usepackage{fancyhdr}
\usepackage{hyperref}
\usepackage{tcolorbox}
\tcbuselibrary{skins, breakable}
\usepackage{fontawesome5}

% =====================================================================
% Shared TikZ styles for the Industrial Security lecture series.
% Loaded by every slide deck and exercise sheet that uses TikZ.
% =====================================================================

\usetikzlibrary{
  arrows.meta,
  shapes,
  shapes.geometric,
  shapes.symbols,
  shapes.multipart,
  positioning,
  fit,
  calc,
  backgrounds,
  decorations.pathmorphing,
  decorations.pathreplacing,
  decorations.text,
  patterns,
  matrix,
  shadows
}

% --- colour palette --------------------------------------------------
\definecolor{isOT}{HTML}{1F4E79}     % OT (deep blue)
\definecolor{isIT}{HTML}{6F2DA8}     % IT (purple)
\definecolor{isSafety}{HTML}{C00000} % safety red
\definecolor{isNet}{HTML}{2E7D32}    % network green
\definecolor{isAttack}{HTML}{B23A48} % attack red
\definecolor{isAccent}{HTML}{E08E0B} % amber
\definecolor{isMute}{HTML}{6B6B6B}   % grey
\definecolor{isLight}{HTML}{EDF2F8}  % light background
\definecolor{isPLC}{HTML}{2C5282}
\definecolor{isHMI}{HTML}{2F855A}
\definecolor{isSCADA}{HTML}{6B46C1}

% --- generic node styles --------------------------------------------
\tikzset{
  isnode/.style={
    draw, rounded corners=2pt, minimum height=8mm, minimum width=18mm,
    font=\small, align=center, thick
  },
  isbox/.style={isnode, fill=isLight},
  plc/.style={isnode, fill=isPLC!15, draw=isPLC, thick},
  hmi/.style={isnode, fill=isHMI!15, draw=isHMI, thick},
  scada/.style={isnode, fill=isSCADA!15, draw=isSCADA, thick},
  sensor/.style={isnode, fill=isNet!10, draw=isNet, minimum height=6mm, minimum width=14mm, font=\scriptsize},
  actuator/.style={isnode, fill=isAccent!15, draw=isAccent, minimum height=6mm, minimum width=14mm, font=\scriptsize},
  firewall/.style={isnode, fill=isAttack!15, draw=isAttack, thick},
  server/.style={isnode, fill=isMute!15, draw=isMute!80, thick},
  switch/.style={isnode, fill=isLight, draw=black!60, minimum height=6mm, minimum width=14mm, font=\scriptsize},
  workstation/.style={isnode, fill=white, draw=isOT, thick},
  attacker/.style={isnode, fill=isAttack!20, draw=isAttack, thick, font=\small\bfseries},
  inetnode/.style={draw, ellipse, minimum width=24mm, minimum height=10mm, font=\scriptsize, fill=isLight, thick},
  process/.style={isnode, fill=isOT!15, draw=isOT, thick, minimum width=24mm},
  zone/.style={
    draw, rounded corners=4pt, thick, dashed, inner sep=4mm
  },
  conduit/.style={-{Stealth[length=2.5mm]}, thick, isNet!80!black},
  attack/.style={-{Stealth[length=2.5mm]}, thick, isAttack, dashed},
  flow/.style={-{Stealth[length=2.5mm]}, thick, isOT!70!black},
  netline/.style={thick, black!60},
  axisline/.style={-{Stealth[length=2mm]}, thick, black!70},
  tag/.style={font=\scriptsize\itshape, fill=white, inner sep=1pt},
  level/.style={
    draw, fill=isLight, minimum width=0.85\linewidth, minimum height=8mm,
    font=\small, anchor=west
  },
  brace/.style={decorate, decoration={brace, amplitude=4pt, raise=2pt}},
  legendbox/.style={draw, rounded corners=2pt, fill=isLight, inner sep=2mm, font=\scriptsize, align=left}
}

% --- handy macros ----------------------------------------------------
% \plcicon, \hmiicon etc as simple labelled boxes; useful inline.
\newcommand{\plcicon}[1][PLC]{\tikz[baseline=-0.6ex]{\node[plc, minimum width=10mm, minimum height=5mm, font=\scriptsize, inner sep=1pt]{#1};}}
\newcommand{\hmiicon}[1][HMI]{\tikz[baseline=-0.6ex]{\node[hmi, minimum width=10mm, minimum height=5mm, font=\scriptsize, inner sep=1pt]{#1};}}
\newcommand{\fwicon}[1][FW]{\tikz[baseline=-0.6ex]{\node[firewall, minimum width=10mm, minimum height=5mm, font=\scriptsize, inner sep=1pt]{#1};}}


\providecommand{\courseAuthor}{Industrial Security Lecture Series}
\providecommand{\courseInstitute}{Department of Computer Science}
\providecommand{\companyLogo}{logo}

\definecolor{slideAccent}{HTML}{00205B}
\definecolor{slideSecondary}{HTML}{00A0DC}

% --- Corporate branding overrides ------------------------------------
\IfFileExists{../../../branding.tex}{% =====================================================================
% Corporate / institutional branding -- single file to customise the
% look of the entire course (slides, notes, exercises, labs).
%
% HOW TO REBRAND IN UNDER A MINUTE:
%   1. Replace `branding/logo.pdf` with your own logo
%      (PDF preferred; PNG and JPG also work -- name it `logo.<ext>`).
%   2. (Optional) change the two colours below to your brand palette.
%   3. (Optional) override the author/institute lines below.
%   4. Re-run `make` at the repo root. That's it.
%
% Everything in this file has sensible defaults; you can delete a line
% to fall back to the built-in defaults, or comment it out.
%
% This file is loaded automatically by the slides, exercise, and lab
% preambles -- you never need to \input it manually.
% =====================================================================

% --- Logo --------------------------------------------------------------
% Path is resolved from each leaf-build directory; the default points at
% the shared `branding/` folder at the repo root. If you put your logo
% somewhere else, set the full or relative path here (without extension):
\renewcommand{\companyLogo}{../../../branding/logo}

% --- Author / institute (appears on the title page footer) ------------
\renewcommand{\courseAuthor}{}
\renewcommand{\courseInstitute}{}

% --- Brand colours ----------------------------------------------------
% slideAccent      = primary brand colour (title bands, accent rule)
% slideSecondary   = secondary brand colour (section badge, highlight)
% Both use HTML hex codes. Keep enough contrast against white text.
\definecolor{slideAccent}{HTML}{00205B}      % deep navy
\definecolor{slideSecondary}{HTML}{00A0DC}   % light blue accent
}{}

\colorlet{labAccent}{slideAccent}
\colorlet{labSec}{slideSecondary}
\definecolor{labCheck}{HTML}{2E7D32}
\definecolor{labWarn}{HTML}{B23A48}

\lstset{
  backgroundcolor=\color{black!4},
  basicstyle=\ttfamily\footnotesize,
  breaklines=true,
  frame=leftline,
  framerule=2pt,
  rulecolor=\color{labAccent},
  numbers=left,
  numberstyle=\tiny\color{black!50},
  showstringspaces=false,
  tabsize=2
}

\setlength{\tabcolsep}{4pt}

% Let TeX add a little extra inter-word stretch as a last resort before
% it reports an Overfull \hbox. Only kicks in on lines that would
% otherwise overflow (long \texttt paths, URLs, CIDRs); never loosens a
% line that already fits. Keeps prose/itemize within the margin.
\setlength{\emergencystretch}{3em}

\pagestyle{fancy}
\fancyhf{}
\renewcommand{\headrulewidth}{0.4pt}
\lhead{\small \courseTitle{} -- Lab \labNumber}
\rhead{\small Maps to Section \labSection}
\cfoot{\thepage}

\newtcolorbox{stepbox}[1]{
  enhanced, breakable, arc=2pt, boxrule=0pt,
  colback=labAccent!7, colframe=labAccent, leftrule=3pt,
  fonttitle=\bfseries\small, coltitle=white,
  colbacktitle=labAccent,
  title={\faTools\ Step #1}
}

\newtcolorbox{warnbox}{
  enhanced, breakable, arc=2pt, boxrule=0pt,
  colback=labWarn!7, colframe=labWarn, leftrule=4pt,
  fonttitle=\bfseries\small, coltitle=white, colbacktitle=labWarn,
  title={\faExclamationTriangle\ Caution}
}

\newtcolorbox{notebox}{
  enhanced, breakable, arc=2pt, boxrule=0pt,
  colback=labAccent!4, colframe=labAccent, leftrule=2pt,
  fonttitle=\bfseries\small, coltitle=white, colbacktitle=labAccent,
  title={\faInfoCircle\ Note}
}

\newtcolorbox{goalbox}{
  enhanced, breakable, arc=2pt, boxrule=0pt,
  colback=labCheck!7, colframe=labCheck, leftrule=3pt,
  fonttitle=\bfseries\small, coltitle=white, colbacktitle=labCheck,
  title={\faBullseye\ Goal}
}

\newtcolorbox{deliverbox}{
  enhanced, breakable, arc=2pt, boxrule=0pt,
  colback=labSec!10, colframe=labSec, leftrule=3pt,
  fonttitle=\bfseries\small, coltitle=white, colbacktitle=labSec,
  title={\faClipboardList\ Deliverable}
}

% --- Solution-sheet boxes (used only by lab-solutions.tex) ----------
\newtcolorbox{solutionbox}[1]{
  enhanced, breakable, arc=2pt, boxrule=0pt,
  colback=labCheck!7, colframe=labCheck, leftrule=3pt,
  fonttitle=\bfseries\small, coltitle=white, colbacktitle=labCheck,
  title={\faCheckCircle\ #1}
}

\newtcolorbox{rubricbox}{
  enhanced, breakable, arc=2pt, boxrule=0pt,
  colback=labSec!7, colframe=labSec, leftrule=3pt,
  fonttitle=\bfseries\small, coltitle=white, colbacktitle=labSec,
  title={\faBalanceScale\ Grading rubric}
}

\titleformat{\section}{\large\bfseries\color{labAccent}}{\thesection}{0.5em}{}

\newcommand{\labheader}{%
  \noindent
  \begin{tikzpicture}
    \fill[labAccent] (0,0) rectangle (\textwidth, -1.6cm);
    \node[anchor=north west, text=white, font=\bfseries\large] at (0.6cm, -0.4cm) {Lab \labNumber{} -- \labTitle};
    \node[anchor=south west, text=white, font=\small] at (0.6cm, -1.3cm) {\courseTitle{} \quad $\bullet$ \quad Maps to Section \labSection{} \quad $\bullet$ \quad \labTime};
    \node[anchor=north east, fill=labSec, text=white, font=\bfseries, inner sep=4pt, rounded corners=2pt]
      at ([xshift=-0.4cm, yshift=-0.4cm]\textwidth,0) {LAB \labNumber};
  \end{tikzpicture}
  \vspace{4mm}
}

% =====================================================================

\usepackage[utf8]{inputenc}
\usepackage[T1]{fontenc}
\usepackage[english]{babel}
\usepackage[a4paper, margin=2cm]{geometry}
\usepackage{graphicx}
\usepackage{listings}
\usepackage{xcolor}
\usepackage{colortbl}
\usepackage{tikz}
\usepackage{amsmath}
\usepackage{amssymb}
\usepackage{booktabs}
\usepackage{enumitem}
\usepackage{titlesec}
\usepackage{fancyhdr}
\usepackage{hyperref}
\usepackage{tcolorbox}
\tcbuselibrary{skins, breakable}
\usepackage{fontawesome5}

% =====================================================================
% Shared TikZ styles for the Industrial Security lecture series.
% Loaded by every slide deck and exercise sheet that uses TikZ.
% =====================================================================

\usetikzlibrary{
  arrows.meta,
  shapes,
  shapes.geometric,
  shapes.symbols,
  shapes.multipart,
  positioning,
  fit,
  calc,
  backgrounds,
  decorations.pathmorphing,
  decorations.pathreplacing,
  decorations.text,
  patterns,
  matrix,
  shadows
}

% --- colour palette --------------------------------------------------
\definecolor{isOT}{HTML}{1F4E79}     % OT (deep blue)
\definecolor{isIT}{HTML}{6F2DA8}     % IT (purple)
\definecolor{isSafety}{HTML}{C00000} % safety red
\definecolor{isNet}{HTML}{2E7D32}    % network green
\definecolor{isAttack}{HTML}{B23A48} % attack red
\definecolor{isAccent}{HTML}{E08E0B} % amber
\definecolor{isMute}{HTML}{6B6B6B}   % grey
\definecolor{isLight}{HTML}{EDF2F8}  % light background
\definecolor{isPLC}{HTML}{2C5282}
\definecolor{isHMI}{HTML}{2F855A}
\definecolor{isSCADA}{HTML}{6B46C1}

% --- generic node styles --------------------------------------------
\tikzset{
  isnode/.style={
    draw, rounded corners=2pt, minimum height=8mm, minimum width=18mm,
    font=\small, align=center, thick
  },
  isbox/.style={isnode, fill=isLight},
  plc/.style={isnode, fill=isPLC!15, draw=isPLC, thick},
  hmi/.style={isnode, fill=isHMI!15, draw=isHMI, thick},
  scada/.style={isnode, fill=isSCADA!15, draw=isSCADA, thick},
  sensor/.style={isnode, fill=isNet!10, draw=isNet, minimum height=6mm, minimum width=14mm, font=\scriptsize},
  actuator/.style={isnode, fill=isAccent!15, draw=isAccent, minimum height=6mm, minimum width=14mm, font=\scriptsize},
  firewall/.style={isnode, fill=isAttack!15, draw=isAttack, thick},
  server/.style={isnode, fill=isMute!15, draw=isMute!80, thick},
  switch/.style={isnode, fill=isLight, draw=black!60, minimum height=6mm, minimum width=14mm, font=\scriptsize},
  workstation/.style={isnode, fill=white, draw=isOT, thick},
  attacker/.style={isnode, fill=isAttack!20, draw=isAttack, thick, font=\small\bfseries},
  inetnode/.style={draw, ellipse, minimum width=24mm, minimum height=10mm, font=\scriptsize, fill=isLight, thick},
  process/.style={isnode, fill=isOT!15, draw=isOT, thick, minimum width=24mm},
  zone/.style={
    draw, rounded corners=4pt, thick, dashed, inner sep=4mm
  },
  conduit/.style={-{Stealth[length=2.5mm]}, thick, isNet!80!black},
  attack/.style={-{Stealth[length=2.5mm]}, thick, isAttack, dashed},
  flow/.style={-{Stealth[length=2.5mm]}, thick, isOT!70!black},
  netline/.style={thick, black!60},
  axisline/.style={-{Stealth[length=2mm]}, thick, black!70},
  tag/.style={font=\scriptsize\itshape, fill=white, inner sep=1pt},
  level/.style={
    draw, fill=isLight, minimum width=0.85\linewidth, minimum height=8mm,
    font=\small, anchor=west
  },
  brace/.style={decorate, decoration={brace, amplitude=4pt, raise=2pt}},
  legendbox/.style={draw, rounded corners=2pt, fill=isLight, inner sep=2mm, font=\scriptsize, align=left}
}

% --- handy macros ----------------------------------------------------
% \plcicon, \hmiicon etc as simple labelled boxes; useful inline.
\newcommand{\plcicon}[1][PLC]{\tikz[baseline=-0.6ex]{\node[plc, minimum width=10mm, minimum height=5mm, font=\scriptsize, inner sep=1pt]{#1};}}
\newcommand{\hmiicon}[1][HMI]{\tikz[baseline=-0.6ex]{\node[hmi, minimum width=10mm, minimum height=5mm, font=\scriptsize, inner sep=1pt]{#1};}}
\newcommand{\fwicon}[1][FW]{\tikz[baseline=-0.6ex]{\node[firewall, minimum width=10mm, minimum height=5mm, font=\scriptsize, inner sep=1pt]{#1};}}


\providecommand{\courseAuthor}{Industrial Security Lecture Series}
\providecommand{\courseInstitute}{Department of Computer Science}
\providecommand{\companyLogo}{logo}

\definecolor{slideAccent}{HTML}{00205B}
\definecolor{slideSecondary}{HTML}{00A0DC}

% --- Corporate branding overrides ------------------------------------
\IfFileExists{../../../branding.tex}{% =====================================================================
% Corporate / institutional branding -- single file to customise the
% look of the entire course (slides, notes, exercises, labs).
%
% HOW TO REBRAND IN UNDER A MINUTE:
%   1. Replace `branding/logo.pdf` with your own logo
%      (PDF preferred; PNG and JPG also work -- name it `logo.<ext>`).
%   2. (Optional) change the two colours below to your brand palette.
%   3. (Optional) override the author/institute lines below.
%   4. Re-run `make` at the repo root. That's it.
%
% Everything in this file has sensible defaults; you can delete a line
% to fall back to the built-in defaults, or comment it out.
%
% This file is loaded automatically by the slides, exercise, and lab
% preambles -- you never need to \input it manually.
% =====================================================================

% --- Logo --------------------------------------------------------------
% Path is resolved from each leaf-build directory; the default points at
% the shared `branding/` folder at the repo root. If you put your logo
% somewhere else, set the full or relative path here (without extension):
\renewcommand{\companyLogo}{../../../branding/logo}

% --- Author / institute (appears on the title page footer) ------------
\renewcommand{\courseAuthor}{}
\renewcommand{\courseInstitute}{}

% --- Brand colours ----------------------------------------------------
% slideAccent      = primary brand colour (title bands, accent rule)
% slideSecondary   = secondary brand colour (section badge, highlight)
% Both use HTML hex codes. Keep enough contrast against white text.
\definecolor{slideAccent}{HTML}{00205B}      % deep navy
\definecolor{slideSecondary}{HTML}{00A0DC}   % light blue accent
}{}

\colorlet{labAccent}{slideAccent}
\colorlet{labSec}{slideSecondary}
\definecolor{labCheck}{HTML}{2E7D32}
\definecolor{labWarn}{HTML}{B23A48}

\lstset{
  backgroundcolor=\color{black!4},
  basicstyle=\ttfamily\footnotesize,
  breaklines=true,
  frame=leftline,
  framerule=2pt,
  rulecolor=\color{labAccent},
  numbers=left,
  numberstyle=\tiny\color{black!50},
  showstringspaces=false,
  tabsize=2
}

\setlength{\tabcolsep}{4pt}

% Let TeX add a little extra inter-word stretch as a last resort before
% it reports an Overfull \hbox. Only kicks in on lines that would
% otherwise overflow (long \texttt paths, URLs, CIDRs); never loosens a
% line that already fits. Keeps prose/itemize within the margin.
\setlength{\emergencystretch}{3em}

\pagestyle{fancy}
\fancyhf{}
\renewcommand{\headrulewidth}{0.4pt}
\lhead{\small \courseTitle{} -- Lab \labNumber}
\rhead{\small Maps to Section \labSection}
\cfoot{\thepage}

\newtcolorbox{stepbox}[1]{
  enhanced, breakable, arc=2pt, boxrule=0pt,
  colback=labAccent!7, colframe=labAccent, leftrule=3pt,
  fonttitle=\bfseries\small, coltitle=white,
  colbacktitle=labAccent,
  title={\faTools\ Step #1}
}

\newtcolorbox{warnbox}{
  enhanced, breakable, arc=2pt, boxrule=0pt,
  colback=labWarn!7, colframe=labWarn, leftrule=4pt,
  fonttitle=\bfseries\small, coltitle=white, colbacktitle=labWarn,
  title={\faExclamationTriangle\ Caution}
}

\newtcolorbox{notebox}{
  enhanced, breakable, arc=2pt, boxrule=0pt,
  colback=labAccent!4, colframe=labAccent, leftrule=2pt,
  fonttitle=\bfseries\small, coltitle=white, colbacktitle=labAccent,
  title={\faInfoCircle\ Note}
}

\newtcolorbox{goalbox}{
  enhanced, breakable, arc=2pt, boxrule=0pt,
  colback=labCheck!7, colframe=labCheck, leftrule=3pt,
  fonttitle=\bfseries\small, coltitle=white, colbacktitle=labCheck,
  title={\faBullseye\ Goal}
}

\newtcolorbox{deliverbox}{
  enhanced, breakable, arc=2pt, boxrule=0pt,
  colback=labSec!10, colframe=labSec, leftrule=3pt,
  fonttitle=\bfseries\small, coltitle=white, colbacktitle=labSec,
  title={\faClipboardList\ Deliverable}
}

% --- Solution-sheet boxes (used only by lab-solutions.tex) ----------
\newtcolorbox{solutionbox}[1]{
  enhanced, breakable, arc=2pt, boxrule=0pt,
  colback=labCheck!7, colframe=labCheck, leftrule=3pt,
  fonttitle=\bfseries\small, coltitle=white, colbacktitle=labCheck,
  title={\faCheckCircle\ #1}
}

\newtcolorbox{rubricbox}{
  enhanced, breakable, arc=2pt, boxrule=0pt,
  colback=labSec!7, colframe=labSec, leftrule=3pt,
  fonttitle=\bfseries\small, coltitle=white, colbacktitle=labSec,
  title={\faBalanceScale\ Grading rubric}
}

\titleformat{\section}{\large\bfseries\color{labAccent}}{\thesection}{0.5em}{}

\newcommand{\labheader}{%
  \noindent
  \begin{tikzpicture}
    \fill[labAccent] (0,0) rectangle (\textwidth, -1.6cm);
    \node[anchor=north west, text=white, font=\bfseries\large] at (0.6cm, -0.4cm) {Lab \labNumber{} -- \labTitle};
    \node[anchor=south west, text=white, font=\small] at (0.6cm, -1.3cm) {\courseTitle{} \quad $\bullet$ \quad Maps to Section \labSection{} \quad $\bullet$ \quad \labTime};
    \node[anchor=north east, fill=labSec, text=white, font=\bfseries, inner sep=4pt, rounded corners=2pt]
      at ([xshift=-0.4cm, yshift=-0.4cm]\textwidth,0) {LAB \labNumber};
  \end{tikzpicture}
  \vspace{4mm}
}


\begin{document}
\labheader

\begin{goalbox}
Stand up a passive ICS-aware IDS (Zeek) on the lab network, baseline the ``normal'' Modbus chatter coming out of OpenPLC, then write and tune a custom detection rule against a simulated attack. By the end you will be able to read a Zeek policy script the way you read a Suricata signature -- and explain why ``policy-as-code'' is the right vocabulary for what we are doing.
\end{goalbox}

\begin{notebox}{\faInfoCircle\ How to read this lab if you come from IT}
You do not need a PLC background. Use these one-line analogues whenever a piece of OT jargon appears:
\begin{itemize}\setlength{\itemsep}{1pt}
  \item \textbf{Modbus TCP} $\approx$ HTTP for the plant -- a request/response protocol on port 502, but with no auth, no TLS, no host header.
  \item \textbf{\texttt{modbus.log}} $\approx$ \texttt{access.log} for the PLC -- one line per request/response, with a verb (\emph{function code}) and an object (register or coil address).
  \item \textbf{Function code} $\approx$ HTTP verb. FC1/2/3/4 = ``GET'' (read). FC5/6/15/16 = ``POST/PUT'' (write). The interesting ones are the writes.
  \item \textbf{\texttt{local.zeek}} $\approx$ a Suricata rules file -- except it is a real (Turing-complete) scripting language, not a flat signature list. We call this \emph{policy-as-code}.
  \item \textbf{\texttt{notice.log}} $\approx$ \texttt{alerts.json} -- one line per Zeek-raised alert, ready for SIEM ingest.
\end{itemize}
\end{notebox}

% ---------------------------------------------------------------------------
\section*{Pre-flight}
% ---------------------------------------------------------------------------

\begin{stepbox}{0 -- Stack up and Zeek alive}
From the repo root, bring up the full lab stack and confirm Zeek is producing logs. The Zeek container shares OpenPLC's network namespace (\texttt{network\_mode: service:openplc} in \texttt{docker-compose.yml}), so its \texttt{eth0} sees every packet to/from the PLC -- the Docker equivalent of a SPAN port hanging off the process switch.
\begin{lstlisting}[language=bash]
cd bachelor/labs/_stack
docker compose up -d
docker compose ps zeek openplc student
\end{lstlisting}
All three services must show \texttt{Up}. Confirm Zeek has begun writing to the host-mounted log directory:
\begin{lstlisting}[language=bash]
ls zeek-logs/
\end{lstlisting}
You should see at least \texttt{conn.log} and \texttt{packet\_filter.log}. \texttt{modbus.log} appears the first time a Modbus packet is seen.
\end{stepbox}

\begin{stepbox}{0a -- Sanity check the dissector}
\begin{lstlisting}[language=bash]
docker compose exec student python3 /labs/scripts/modbus_read.py
sleep 2
ls zeek-logs/ | grep modbus
tail -n 2 zeek-logs/modbus.log
\end{lstlisting}
Expected: two JSON lines, one with \texttt{"pdu\_type":"REQ"} and one with \texttt{"pdu\_type":"RESP"}, both naming the function (e.g.\ \texttt{READ\_HOLDING\_REGISTERS}). If you see no \texttt{modbus.log}, Zeek either did not bind to \texttt{eth0} or the Modbus analyzer is not loaded -- check \texttt{docker compose logs zeek}.
\end{stepbox}

% ---------------------------------------------------------------------------
\section*{Part A -- Baseline}
% ---------------------------------------------------------------------------

\begin{stepbox}{1 -- Generate one minute of legitimate traffic}
Before you can detect ``abnormal'' you need to know what ``normal'' looks like. Start the polling helper in the background and let it run for sixty seconds:
\begin{lstlisting}[language=bash]
docker compose exec -d student python3 /labs/scripts/modbus_poll_loop.py
sleep 60
docker compose exec student pkill -f modbus_poll_loop
\end{lstlisting}
The helper issues one \texttt{READ\_HOLDING\_REGISTERS} (FC3) every two seconds -- the kind of polling rhythm an HMI maintains in real life.
\end{stepbox}

\begin{stepbox}{2 -- Inspect \texttt{conn.log} (the netflow view)}
Look at the connection summary first. \texttt{conn.log} is Zeek's equivalent of a NetFlow record: one line per TCP/UDP flow, with byte counts, duration, and the application-layer service Zeek decoded:
\begin{lstlisting}[language=bash]
tail -n 5 zeek-logs/conn.log | jq '{ts, "src":."id.orig_h", "dst":."id.resp_h",
                                     dport:."id.resp_p", service, duration,
                                     orig_bytes, resp_bytes}'
\end{lstlisting}
Expected: at least one row with \texttt{"service":"modbus"}, source \texttt{172.28.0.20} (the student container), destination \texttt{172.28.0.10} (OpenPLC), destination port \texttt{502}. \emph{This is the only flow your baseline should show on port 502.} Anything else is a finding.
\end{stepbox}

\begin{stepbox}{3 -- Inspect \texttt{modbus.log} (the application view)}
Now drop down one layer to the Modbus dissector output:
\begin{lstlisting}[language=bash]
jq -r '[.ts, ."id.orig_h", ."id.resp_h", .func, .pdu_type] | @tsv' \
    zeek-logs/modbus.log | head -10
\end{lstlisting}
Expected output (timestamps will differ):
\begin{lstlisting}[language=bash]
1779125580.566   172.28.0.20   172.28.0.10   READ_HOLDING_REGISTERS   REQ
1779125580.570   172.28.0.10   172.28.0.20   READ_HOLDING_REGISTERS   RESP
1779125582.567   172.28.0.20   172.28.0.10   READ_HOLDING_REGISTERS   REQ
1779125582.571   172.28.0.10   172.28.0.20   READ_HOLDING_REGISTERS   RESP
...
\end{lstlisting}
\textbf{Record the baseline as a set of \emph{(source IP, destination IP, function)} triples} -- this is the inventory of ``legitimate Modbus conversations'' for the rest of the lab:
\begin{lstlisting}[language=bash]
jq -r '"\(."id.orig_h")\t\(."id.resp_h")\t\(.func)"' zeek-logs/modbus.log \
    | sort -u > /tmp/baseline_triples.tsv
cat /tmp/baseline_triples.tsv
\end{lstlisting}
You should see exactly one read triple and the matching response. Keep this file -- you will compare against it after the attack.
\end{stepbox}

% ---------------------------------------------------------------------------
\section*{Part B -- Read the existing detection policy}
% ---------------------------------------------------------------------------

\begin{stepbox}{4 -- Walk through \texttt{local.zeek} in plain English}
Open \texttt{bachelor/labs/\_stack/zeek/local.zeek} in your editor. The file is short on purpose -- read every line. Below is the same script annotated for an IT reader:
\begin{lstlisting}[language=bash]
@load base/protocols/modbus                       # turn on the dissector
@load policy/protocols/modbus/known-masters-slaves# enumerate hosts seen on 502

redef LogAscii::use_json = T;                     # log as JSON, not TSV

module ModbusWatch;                               # our own namespace
export {
    redef enum Notice::Type += { Unsolicited_Write };
}

event modbus_message(c: connection, headers: ModbusHeaders, is_orig: bool) {
    if (!is_orig) return;                         # only inspect requests
    if (headers$function_code in set(5, 6, 15, 16)) {
        NOTICE([$note=Unsolicited_Write,
                $msg=fmt("Modbus write from %s to %s (FC=%d)",
                         c$id$orig_h, c$id$resp_h, headers$function_code),
                $conn=c]);
    }
}
\end{lstlisting}
What each line does, in IT terms:
\begin{itemize}\setlength{\itemsep}{1pt}
  \item \texttt{@load base/protocols/modbus} -- enable the Modbus \emph{analyzer} (``install the Wireshark dissector inside Zeek''). Without it, port 502 is just a TCP byte stream.
  \item \texttt{@load policy/protocols/modbus/known-masters-slaves} -- a built-in policy that produces \texttt{known\_modbus.log}, a unique-pair inventory of who talked Modbus to whom. Asset-discovery freebie.
  \item \texttt{redef LogAscii::use\_json = T} -- emit JSON so SIEMs ingest it without a custom parser.
  \item \texttt{module ModbusWatch} -- declare a namespace so our notice type does not clash with Zeek-shipped ones.
  \item \texttt{redef enum Notice::Type += \{ Unsolicited\_Write \}} -- register a new alert category. Notices are Zeek's first-class alert primitive; they land in \texttt{notice.log} with a stable identifier you can pivot on in the SOC.
  \item \texttt{event modbus\_message(\dots)} -- Zeek calls this handler every time the dissector finishes parsing a Modbus PDU. \texttt{is\_orig = T} means ``client to server'' (a request).
  \item \texttt{set(5, 6, 15, 16)} is the write-side function codes: FC5 = write single coil, FC6 = write single register, FC15 = write multiple coils, FC16 = write multiple registers.
  \item \texttt{NOTICE([\dots])} raises an alert. The \texttt{\$conn=c} field attaches the full connection 5-tuple so the SOC can pivot to \texttt{conn.log}.
\end{itemize}
\textbf{The conceptual point.} This file is a tiny Python-flavoured DSL doing exactly what your Suricata or Snort rules file would do -- the difference is that here we have full programmatic state (sets, tables, loops, timers, file I/O). That is why people call Zeek ``policy-as-code'' instead of ``signature matching.''
\end{stepbox}

% ---------------------------------------------------------------------------
\section*{Part C -- Trigger the existing detection}
% ---------------------------------------------------------------------------

\begin{stepbox}{5 -- Fire the attack and watch \texttt{notice.log}}
Re-run a few baseline polls to make sure the noise is still flowing, then inject a single unauthenticated write to coil~0:
\begin{lstlisting}[language=bash]
docker compose exec -d student python3 /labs/scripts/modbus_poll_loop.py
sleep 5
docker compose exec student python3 /labs/scripts/modbus_inject.py
sleep 2
tail -n 3 zeek-logs/notice.log
\end{lstlisting}
Expected line (JSON, one field per attribute):
\begin{lstlisting}[language=bash]
{"ts":...,"note":"ModbusWatch::Unsolicited_Write",
 "msg":"Modbus write from 172.28.0.20 to 172.28.0.10 (FC=5)",
 "id.orig_h":"172.28.0.20","id.resp_h":"172.28.0.10",...}
\end{lstlisting}
Stop the poller before continuing:
\begin{lstlisting}[language=bash]
docker compose exec student pkill -f modbus_poll_loop
\end{lstlisting}
\end{stepbox}

\begin{stepbox}{6 -- Confirm the alert in \texttt{modbus.log} too}
Cross-check by grepping the application log for any write function code. There should be exactly one matching record:
\begin{lstlisting}[language=bash]
jq -r 'select(.func | test("WRITE")) | [.ts, ."id.orig_h", .func] | @tsv' \
    zeek-logs/modbus.log
\end{lstlisting}
Expected: one line per injected write, naming \texttt{WRITE\_SINGLE\_COIL}. If \texttt{notice.log} has the alert but \texttt{modbus.log} does not show the write, your dissector is misconfigured -- not your detection.
\end{stepbox}

% ---------------------------------------------------------------------------
\section*{Part D -- Write your own detection}
% ---------------------------------------------------------------------------

\begin{stepbox}{7 -- The new rule: ``a source that has never written before''}
\textbf{Problem statement.} The default rule alerts on \emph{every} write. In a real plant the engineering workstation legitimately writes during change windows -- so the rule is noisy. A better question is: \emph{has this source IP ever written before?} The first write from a previously read-only source is the detection that catches Stuxnet-shaped attacks.

\textbf{Approach.} Maintain a stateful set of \texttt{(src,dst)} pairs that have ever sent a write. On each write packet, check the set; if the pair is new, raise a different notice (\texttt{First\_Time\_Writer}) and add the pair.

Open \texttt{bachelor/labs/\_stack/zeek/local.zeek} and \textbf{append} the following block (do not delete the existing handler -- the two will coexist):
\begin{lstlisting}[language=bash]
# --- Rule 2: first-time writer ----------------------------------------
module ModbusWatch;

export {
    redef enum Notice::Type += { First_Time_Writer };
}

# (src,dst) pairs we have ever seen issuing a write
global seen_writers: set[addr, addr] = set();

event modbus_message(c: connection, headers: ModbusHeaders, is_orig: bool)
    &priority=-5      # run after the existing handler
{
    if (!is_orig) return;
    if (headers$function_code !in set(5, 6, 15, 16)) return;

    local src = c$id$orig_h;
    local dst = c$id$resp_h;

    if ([src, dst] in seen_writers) return;     # already known: be quiet
    add seen_writers[src, dst];

    NOTICE([$note=First_Time_Writer,
            $msg=fmt("First Modbus write ever seen from %s to %s (FC=%d)",
                     src, dst, headers$function_code),
            $conn=c]);
}
\end{lstlisting}
Things worth pointing out:
\begin{itemize}\setlength{\itemsep}{1pt}
  \item \texttt{global seen\_writers: set[addr, addr]} -- Zeek has first-class sets keyed on any tuple of typed values. This is the bit you cannot do in Suricata.
  \item \texttt{\&priority=-5} -- when two handlers respond to the same event, Zeek runs the higher priority first. Negative priority means ``after the other one,'' so Rule~1's alert is chronologically before Rule~2's in \texttt{notice.log}.
  \item State is per Zeek process, in memory. It resets every time you restart the container. For real deployments you would persist via \texttt{Broker} or an external store -- out of scope here.
\end{itemize}
\end{stepbox}

\begin{stepbox}{8 -- Reload Zeek and retrigger the attack}
\begin{lstlisting}[language=bash]
docker compose restart zeek
sleep 3
docker compose logs --tail=20 zeek    # confirm clean startup, no syntax errors
\end{lstlisting}
If Zeek fails to start, the logs show the exact line and column of the syntax error -- fix and try again.

Generate fresh baseline, then attack:
\begin{lstlisting}[language=bash]
docker compose exec -d student python3 /labs/scripts/modbus_poll_loop.py
sleep 5
docker compose exec student python3 /labs/scripts/modbus_inject.py
sleep 2
docker compose exec student pkill -f modbus_poll_loop
\end{lstlisting}
Inspect \texttt{notice.log}:
\begin{lstlisting}[language=bash]
jq -r 'select(.note | test("ModbusWatch")) | [.ts, .note, .msg] | @tsv' \
    zeek-logs/notice.log
\end{lstlisting}
Expected: \textbf{two} ModbusWatch notices in chronological order -- first \texttt{Unsolicited\_Write} (Rule~1), then \texttt{First\_Time\_Writer} (Rule~2). If you run \texttt{modbus\_inject.py} a second time without restarting Zeek, only Rule~1 fires -- Rule~2 stays silent because the writer is no longer first-time. \emph{Confirm this experimentally.}
\end{stepbox}

% ---------------------------------------------------------------------------
\section*{Part E -- Tune one false positive}
% ---------------------------------------------------------------------------

\begin{stepbox}{9 -- Pick a noise source and silence it}
In this lab the student workstation (\texttt{172.28.0.20}) \emph{is} our sanctioned engineering host. Treating its writes as alerts is the textbook false-positive case in any new IDS deployment.

Add a \emph{suppression}: any write where the source is \texttt{172.28.0.20} and the destination is \texttt{172.28.0.10} is acknowledged engineering-workstation behaviour. \textbf{Append} this block below your second rule:
\begin{lstlisting}[language=bash]
# --- Tuning: trusted engineering workstation --------------------------
module ModbusWatch;

# (src,dst) pairs that are sanctioned engineering channels
global trusted_writers: set[addr, addr] = {
    [172.28.0.20, 172.28.0.10],
};

hook Notice::policy(n: Notice::Info)
{
    if (n$note != Unsolicited_Write && n$note != First_Time_Writer) return;
    if (!n?$conn) return;
    if ([n$conn$id$orig_h, n$conn$id$resp_h] in trusted_writers)
        add n$actions[Notice::ACTION_IGNORE];
}
\end{lstlisting}
Why a \texttt{Notice::policy} hook rather than editing the rules themselves? Separation of concerns. Detection logic answers \emph{``is this suspicious?''} The notice policy answers \emph{``do we care right now?''} A real SOC tunes the second without touching the first -- so on day two, when the engineering host moves to \texttt{172.28.0.21}, you update one set.

Reload and verify the noise is gone:
\begin{lstlisting}[language=bash]
docker compose restart zeek
sleep 3
> zeek-logs/notice.log            # truncate for a clean before/after
docker compose exec -d student python3 /labs/scripts/modbus_poll_loop.py
sleep 5
docker compose exec student python3 /labs/scripts/modbus_inject.py
sleep 2
docker compose exec student pkill -f modbus_poll_loop
wc -l zeek-logs/notice.log
\end{lstlisting}
Expected: \texttt{0 zeek-logs/notice.log}. The polls are still logged in \texttt{modbus.log} and the write is still logged in \texttt{modbus.log} -- nothing is hidden, only \emph{the alert} is suppressed. Real-world principle: never delete the evidence, only suppress the page.
\end{stepbox}

% ---------------------------------------------------------------------------
\section*{Part F -- Layered defence diagram and 62443 mapping}
% ---------------------------------------------------------------------------

\begin{stepbox}{10 -- Draw the layered defence for this stack}
On paper (or in any diagramming tool), draw the pipeline below left-to-right, naming the host or container at each stage:
\begin{center}
\begin{tikzpicture}[node distance=4mm, every node/.style={font=\scriptsize}]
  \node[draw, fill=labAccent!10, minimum height=10mm, minimum width=22mm, align=center] (a) {Wire / netns\\(\texttt{eth0})};
  \node[draw, fill=labAccent!10, minimum height=10mm, minimum width=22mm, align=center, right=of a] (b) {Zeek\\(dissector)};
  \node[draw, fill=labAccent!10, minimum height=10mm, minimum width=22mm, align=center, right=of b] (c) {\texttt{local.zeek}\\(policy)};
  \node[draw, fill=labAccent!10, minimum height=10mm, minimum width=22mm, align=center, right=of c] (d) {\texttt{notice.log}\\(JSON)};
  \node[draw, fill=labAccent!10, minimum height=10mm, minimum width=22mm, align=center, right=of d] (e) {SIEM\\forward};
  \node[draw, fill=labAccent!10, minimum height=10mm, minimum width=22mm, align=center, right=of e] (f) {SOC\\playbook};
  \draw[->, thick] (a) -- (b);
  \draw[->, thick] (b) -- (c);
  \draw[->, thick] (c) -- (d);
  \draw[->, thick] (d) -- (e);
  \draw[->, thick] (e) -- (f);
\end{tikzpicture}
\end{center}
Label each stage with the IEC~62443-3-3 Foundational Requirement it primarily satisfies (same FR vocabulary as the slides). Write the FR number under each box in your sketch:

\renewcommand{\arraystretch}{1.15}
\begin{center}
\begin{tabular}{p{26mm}p{18mm}p{78mm}}
\toprule
\textbf{Stage} & \textbf{FR} & \textbf{Why} \\
\midrule
PCAP capture (netns, SPAN, TAP) & FR3 & System integrity: monitoring channel cannot be tampered with by an attacker on the data plane. \\
Zeek dissector            & FR6 & Timely response: turning bytes into structured events is the precondition for any alert. \\
\texttt{local.zeek} policy & FR6 & Detection content. Without rules, the dissector is just an expensive logger. \\
\texttt{notice.log}        & FR6 & Centralised logging. \\
SIEM forwarding            & FR5/FR6 & Restricted data flow (one-way push out of the OT zone) plus timely response. \\
SOC playbook               & FR6 & Human-in-the-loop. The playbook is where ``alert'' becomes ``action.'' \\
\bottomrule
\end{tabular}
\end{center}

Reflection question (write the answer in your notes): \emph{which one stage, if it failed silently, would defeat all the others?} Argue for one answer in two sentences.
\end{stepbox}

% ---------------------------------------------------------------------------
\section*{Part G -- Stretch: one-line incident summary}
% ---------------------------------------------------------------------------

\begin{stepbox}{11 -- Build an incident one-liner (stretch goal)}
The classic Zeek tool is \texttt{zeek-cut}, which selects columns from a TSV log. Our \texttt{local.zeek} sets \texttt{LogAscii::use\_json = T}, so \texttt{notice.log} is JSON -- the modern equivalent is \texttt{jq}. Both paths are shown.

\textbf{Path A: \texttt{jq} on the current JSON logs} (recommended):
\begin{lstlisting}[language=bash]
jq -r 'select(.note=="ModbusWatch::First_Time_Writer") |
       "[\(.ts | todate)] src=\(."id.orig_h") fc=write attacker_msg=\(.msg)"' \
    zeek-logs/notice.log
\end{lstlisting}
Correlate with the write payload from \texttt{modbus.log} (which carries the register address and value) and join on the connection UID:
\begin{lstlisting}[language=bash]
UID=$(jq -r 'select(.note=="ModbusWatch::First_Time_Writer") | .uid' \
        zeek-logs/notice.log | tail -n 1)
jq -r --arg u "$UID" \
   'select(.uid==$u and .pdu_type=="REQ") |
    "alert ts=\(.ts) src=\(."id.orig_h") dst=\(."id.resp_h") func=\(.func)
                  addr=\(.address // "n/a") value=\(.values // "n/a")"' \
   zeek-logs/modbus.log
\end{lstlisting}
\textbf{Path B: classic \texttt{zeek-cut}.} Toggle Zeek back to TSV first by commenting out the \texttt{redef LogAscii::use\_json = T;} line in \texttt{local.zeek}, restart Zeek, retrigger the attack, then:
\begin{lstlisting}[language=bash]
docker compose exec zeek bash -c \
   'cat /usr/local/zeek/logs/notice.log | \
    zeek-cut -d ts src note msg | grep First_Time_Writer'
\end{lstlisting}
Expected output (single line, columns tab-separated):
\begin{lstlisting}[language=bash]
2026-05-19T14:32:07+0000   172.28.0.20   ModbusWatch::First_Time_Writer  ...
\end{lstlisting}
Either path gives the SOC the four facts they need to triage in under a minute: \emph{when, who, what, where}. Paste your one-liner into the deliverable.
\end{stepbox}

% ---------------------------------------------------------------------------
\section*{Wrap-up}
% ---------------------------------------------------------------------------

\begin{deliverbox}
Hand in a single PDF or text file with:
\begin{enumerate}\setlength{\itemsep}{1pt}
  \item Your \texttt{/tmp/baseline\_triples.tsv} (the legitimate Modbus inventory).
  \item The first 5 lines of \texttt{modbus.log} during baseline and the offending lines during the attack.
  \item \texttt{notice.log} excerpts at three checkpoints: (i) after Step~5 with one alert, (ii) after Step~8 with two alerts, (iii) after Step~9 with zero alerts.
  \item Your modified \texttt{local.zeek}, complete.
  \item Your layered-defence sketch with FR labels (Step~10).
  \item The one-line incident summary from Step~11.
  \item A 5-line answer to: \emph{``If this alert lands in the SOC at 03:00, who is paged, what is their first action, and what is the one piece of evidence they should preserve before doing anything else?''}
\end{enumerate}
\end{deliverbox}

\begin{notebox}{\faLightbulb\ What you should have learned}
\begin{itemize}\setlength{\itemsep}{1pt}
  \item A passive ICS IDS is a dissector plus a policy script. The dissector turns wire bytes into structured events; the policy decides what is notable.
  \item \texttt{local.zeek} is \emph{code}, not a config file -- stateful sets and event handlers let you express detections that Suricata cannot.
  \item Tuning is a separate concern from detection. Use \texttt{Notice::policy} hooks to silence noise without weakening the rule.
  \item Every layer in the pipeline (capture, dissect, decide, log, forward, respond) maps onto an IEC~62443-3-3 FR -- and a single silent failure anywhere in the chain breaks the whole detection.
\end{itemize}
\end{notebox}

\begin{warnbox}
Everything in this lab runs against containers you own on a private Docker bridge. Running Zeek on a production network is fine -- it is passive by definition -- but you must use a real SPAN port or a hardware TAP, not a shared network namespace, and you must have written authorisation from the network owner. Writing detection rules against live OT traffic without an agreed change window is also out of scope: a buggy rule that emits a NOTICE on every packet can fill a disk and crash the analyser.
\end{warnbox}

\end{document}
